\documentclass[11pt,spanish]{article}

% =========================================================
% PLANTILLA BASE PARA INFORMES ACADÉMICOS / TÉCNICOS
% =========================================================
% Instrucciones rápidas:
% 1. Edite los datos de la sección "DATOS DEL DOCUMENTO".
% 2. Reemplace los textos marcados como TODO.
% 3. Use las tablas de ejemplo como base para nuevas tablas.
% 4. Para tablas anchas, use el entorno widepage incluido abajo.
% 5. Las imágenes se cargan con \SafeImage, que no falla si el archivo no existe.

% --- Idioma y codificación ---
\usepackage[spanish,es-tabla]{babel}
\usepackage[T1]{fontenc}
\usepackage[utf8]{inputenc}

% --- Márgenes / papel ---
\usepackage{geometry}
\geometry{a4paper, margin=2.7cm}

% --- Matemáticas ---
\usepackage{amsmath}

% --- Tablas, enlaces y elementos visuales ---
\usepackage{array}
\usepackage{tabularx}
\usepackage{booktabs}
\usepackage{longtable}
\usepackage{graphicx}
\usepackage{float}
\usepackage{xcolor}
\usepackage{tikz}
\usetikzlibrary{positioning,shapes.geometric,arrows.meta,calc}
\usepackage{enumitem}
\usepackage{ragged2e}
\usepackage[hidelinks]{hyperref}

% --- Encabezado y pie de página ---
\usepackage[myheadings]{fancyhdr}
\usepackage{lastpage}
\setlength{\headheight}{30pt}
\setlength{\footskip}{30pt}

% --- Interlineado ---
\linespread{1.2}
\sloppy

% =========================================================
% DATOS DEL DOCUMENTO
% =========================================================
\newcommand{\Institucion}{Nombre de la institución}
\newcommand{\Campus}{Nombre del campus o sede}
\newcommand{\DocumentoTitulo}{Título del documento}
\newcommand{\DocumentoSubtitulo}{Subtítulo opcional del documento}
\newcommand{\Curso}{Nombre del curso o proyecto}
\newcommand{\Profesor}{Nombre del profesor o responsable}
\newcommand{\FechaDocumento}{Día de mes de año}
\newcommand{\PeriodoAcademico}{Periodo académico}
\newcommand{\AnioDocumento}{Año}

% Autores. Edite, agregue o elimine líneas según corresponda.
\newcommand{\AutoresPortada}{%
    Nombre Apellido: Identificación\\
    Nombre Apellido: Identificación\\
    Nombre Apellido: Identificación%
}

% Datos que aparecen en el encabezado.
\newcommand{\DocHeaderLeft}{Proyecto / Curso}
\newcommand{\DocHeaderRight}{Nombre corto del informe}
\newcommand{\DocFooterRight}{Página~\thepage\ de~\pageref{LastPage}}

% Metadatos PDF.
\title{\DocumentoTitulo}
\author{\AutoresPortada}
\date{\FechaDocumento}

% =========================================================
% CONFIGURACIÓN DE TABLAS Y UTILIDADES
% =========================================================
\newcolumntype{Y}{>{\RaggedRight\arraybackslash}X}
\renewcommand{\arraystretch}{1.22}
\setlength{\LTpre}{0.5em}
\setlength{\LTpost}{0.5em}

% Imagen segura: si el archivo no existe, muestra un recuadro de marcador.
\newcommand{\SafeImage}[2][0.92\textwidth]{%
    \IfFileExists{#2}{%
        \includegraphics[width=#1]{#2}%
    }{%
        \fbox{%
            \begin{minipage}[c][4cm][c]{#1}
            \centering
            Imagen pendiente:\\
            \texttt{#2}
            \end{minipage}%
        }%
    }%
}

% Figura reutilizable.
\newcommand{\EvidenceFigure}[3]{%
    \begin{figure}[H]
    \centering
    \SafeImage{#1}
    \caption{#2}
    \label{#3}
    \end{figure}%
}

% =========================================================
% ENCABEZADO Y PIE
% =========================================================
\pagestyle{fancy}
\fancyhf{}
\fancyhead[L]{\DocHeaderLeft}
\fancyhead[R]{\DocHeaderRight}
\renewcommand{\headrulewidth}{0.4pt}
\renewcommand{\footrulewidth}{0.4pt}
\fancyfoot[R]{\DocFooterRight}

% Estilo equivalente para páginas horizontales.
\fancypagestyle{widefancy}{%
    \fancyhf{}%
    \fancyhead[L]{\DocHeaderLeft}%
    \fancyhead[R]{\DocHeaderRight}%
    \renewcommand{\headrulewidth}{0.4pt}%
    \renewcommand{\footrulewidth}{0.4pt}%
    \fancyfoot[R]{\DocFooterRight}%
}

% =========================================================
% PÁGINAS HORIZONTALES PARA TABLAS ANCHAS
% =========================================================
\makeatletter
\newcommand{\SyncPageBuilderDimensions}{%
    \global\vsize=\textheight
    \global\@colroom=\textheight
    \global\@colht=\textheight
}
\makeatother

\newenvironment{widepage}{%
    \clearpage
    \hypersetup{pageanchor=false}%
    \paperwidth=297mm
    \paperheight=210mm
    \pdfpagewidth=297mm
    \pdfpageheight=210mm
    \newgeometry{left=2.7cm,right=2.7cm,top=2.7cm,bottom=2.7cm}%
    \setlength{\textwidth}{243mm}%
    \setlength{\textheight}{156mm}%
    \setlength{\linewidth}{\textwidth}%
    \setlength{\columnwidth}{\textwidth}%
    \hsize=\textwidth
    \setlength{\headwidth}{\textwidth}%
    \SyncPageBuilderDimensions
    \pagestyle{widefancy}%
}{%
    \clearpage
    \paperwidth=210mm
    \paperheight=297mm
    \pdfpagewidth=210mm
    \pdfpageheight=297mm
    \restoregeometry
    \setlength{\textwidth}{156mm}%
    \setlength{\textheight}{243mm}%
    \setlength{\linewidth}{\textwidth}%
    \setlength{\columnwidth}{\textwidth}%
    \hsize=\textwidth
    \setlength{\headwidth}{\textwidth}%
    \SyncPageBuilderDimensions
    \hypersetup{pageanchor=true}%
    \pagestyle{fancy}%
}

% =========================================================
% PORTADA
% =========================================================
\newcommand{\MakeCover}{%
\begin{titlepage}
\thispagestyle{empty}
\centering

{\bfseries \huge \Institucion \par}
\vspace{0.25cm}
{\LARGE \Campus \par}

\vfill

{\huge \DocumentoTitulo \par}
\vspace{0.4cm}
{\Large \DocumentoSubtitulo \par}

\vfill

{\bfseries \LARGE Profesor(a) / Responsable: \par}
{\Large \Profesor \par}

\vfill

{\bfseries \LARGE Estudiantes / Autores: \par}
{\Large \AutoresPortada \par}

\vfill

{\bfseries\LARGE Fecha: \par}
{\Large \FechaDocumento \par}

\vfill

{\LARGE \PeriodoAcademico \par}
{\LARGE \AnioDocumento \par}

\end{titlepage}%
}

% =========================================================
% DOCUMENTO
% =========================================================
\begin{document}
\hypersetup{pageanchor=false}

% =========================
% Portada
% =========================
\MakeCover

% =========================
% Índice
% =========================
\pagenumbering{roman}
\pagestyle{empty}
\tableofcontents
\clearpage
\pagenumbering{arabic}
\hypersetup{pageanchor=true}
\pagestyle{fancy}

% =========================
% Contenido
% =========================
\section{Introducción}

TODO: Escriba una introducción general del documento, el contexto del trabajo, el alcance evaluado y la finalidad del informe.

\section{Contexto del proyecto}

TODO: Describa el sistema, producto, organización o situación evaluada. Incluya el propósito del proyecto y los módulos o áreas consideradas.

\section{Objetivos}

\subsection{Objetivo general}

TODO: Redacte el objetivo general del documento.

\subsection{Objetivos específicos}

\begin{enumerate}[label=\arabic*.]
    \item TODO: Objetivo específico 1.
    \item TODO: Objetivo específico 2.
    \item TODO: Objetivo específico 3.
\end{enumerate}

\section{Plan de trabajo o plan de calidad}

TODO: Explique la metodología utilizada, responsables, herramientas, criterios de aceptación, métricas o cualquier estrategia aplicada.

\subsection{Alcance y responsabilidades}

\begin{table}[H]
\centering
\caption{Módulos o áreas incluidos en el alcance}
\label{tab:alcance-incluido}
\begin{tabularx}{\textwidth}{p{4.2cm} p{3cm} Y}
\toprule
\textbf{Área / Módulo} & \textbf{Responsable} & \textbf{Actividades principales} \\
\midrule
TODO & TODO & TODO \\
TODO & TODO & TODO \\
TODO & TODO & TODO \\
\bottomrule
\end{tabularx}
\end{table}

\begin{table}[H]
\centering
\caption{Elementos excluidos del alcance}
\label{tab:alcance-excluido}
\begin{tabularx}{\textwidth}{p{4.5cm} Y}
\toprule
\textbf{Elemento excluido} & \textbf{Justificación} \\
\midrule
TODO & TODO \\
TODO & TODO \\
TODO & TODO \\
\bottomrule
\end{tabularx}
\end{table}

\subsection{Requisitos, criterios o elementos evaluados}

\begin{table}[H]
\centering
\caption{Requisitos o criterios evaluados}
\label{tab:requisitos}
\begin{tabularx}{\textwidth}{p{1.8cm} p{2.6cm} Y p{2.7cm}}
\toprule
\textbf{ID} & \textbf{Tipo} & \textbf{Descripción} & \textbf{Responsable} \\
\midrule
REQ-01 & Funcional & TODO & TODO \\
REQ-02 & No funcional & TODO & TODO \\
REQ-03 & Funcional & TODO & TODO \\
REQ-04 & No funcional & TODO & TODO \\
\bottomrule
\end{tabularx}
\end{table}

\subsection{Herramientas utilizadas}

\begin{table}[H]
\centering
\caption{Herramientas utilizadas}
\label{tab:herramientas}
\begin{tabularx}{\textwidth}{p{3.5cm} Y}
\toprule
\textbf{Herramienta} & \textbf{Uso dentro del proyecto} \\
\midrule
TODO & TODO \\
TODO & TODO \\
TODO & TODO \\
\bottomrule
\end{tabularx}
\end{table}

\subsection{Métricas y criterios de aceptación}

\begin{table}[H]
\centering
\caption{Métricas y criterios de aceptación}
\label{tab:metricas-aceptacion}
\begin{tabularx}{\textwidth}{p{4.1cm} p{3.2cm} p{3.2cm} Y}
\toprule
\textbf{Métrica} & \textbf{Aceptación total} & \textbf{Aceptación parcial} & \textbf{Rechazo} \\
\midrule
TODO & TODO & TODO & TODO \\
TODO & TODO & TODO & TODO \\
TODO & TODO & TODO & TODO \\
\bottomrule
\end{tabularx}
\end{table}

\section{Matriz de trazabilidad}

TODO: Explique cómo se relacionan los requisitos o criterios con los casos, evidencias, resultados o responsables. La siguiente tabla está configurada en página horizontal para permitir más columnas.

\begin{widepage}
{\small
\setlength{\tabcolsep}{3pt}
\begin{longtable}{p{1.5cm} p{5.3cm} p{4.4cm} p{3.4cm} p{1.9cm} p{2.2cm} p{2.1cm}}
\caption{Matriz de trazabilidad}\label{tab:trazabilidad}\\
\toprule
\textbf{Req.} & \textbf{Descripción} & \textbf{Casos / Evidencias} & \textbf{Tipo} & \textbf{Resultado} & \textbf{Defecto / Hallazgo} & \textbf{Resp.} \\
\midrule
\endfirsthead
\toprule
\textbf{Req.} & \textbf{Descripción} & \textbf{Casos / Evidencias} & \textbf{Tipo} & \textbf{Resultado} & \textbf{Defecto / Hallazgo} & \textbf{Resp.} \\
\midrule
\endhead
REQ-01 & TODO & TODO & TODO & TODO & TODO & TODO \\
REQ-02 & TODO & TODO & TODO & TODO & TODO & TODO \\
REQ-03 & TODO & TODO & TODO & TODO & TODO & TODO \\
\bottomrule
\end{longtable}
}
\end{widepage}

\section{Casos de prueba, actividades o evidencias ejecutadas}

\subsection{Bloque 1}

TODO: Describa los casos, actividades o pruebas realizadas en este bloque.

\begin{table}[H]
\centering
\caption{Resumen de resultados del bloque 1}
\label{tab:resultados-bloque-1}
\begin{tabularx}{\textwidth}{p{2.8cm} Y p{3cm} p{3cm}}
\toprule
\textbf{ID} & \textbf{Descripción} & \textbf{Resultado esperado} & \textbf{Resultado obtenido} \\
\midrule
CP-001 & TODO & TODO & TODO \\
CP-002 & TODO & TODO & TODO \\
CP-003 & TODO & TODO & TODO \\
\bottomrule
\end{tabularx}
\end{table}

\EvidenceFigure{images/evidencia-ejemplo.png}{Evidencia visual de ejemplo}{fig:evidencia-ejemplo}

\subsection{Bloque 2}

TODO: Describa un segundo conjunto de actividades, pruebas o resultados.

\begin{table}[H]
\centering
\caption{Resultados cuantitativos de ejemplo}
\label{tab:resultados-cuantitativos}
\begin{tabularx}{\textwidth}{p{4cm} p{3cm} p{3cm} Y}
\toprule
\textbf{Indicador} & \textbf{Valor esperado} & \textbf{Valor obtenido} & \textbf{Interpretación} \\
\midrule
TODO & TODO & TODO & TODO \\
TODO & TODO & TODO & TODO \\
TODO & TODO & TODO & TODO \\
\bottomrule
\end{tabularx}
\end{table}

\subsection{Resumen consolidado}

\begin{widepage}
{\small
\setlength{\tabcolsep}{3pt}
\begin{longtable}{p{2.2cm} p{4.5cm} p{2.6cm} p{3.4cm} p{2.4cm} p{6.2cm}}
\caption{Resumen consolidado de casos, actividades o evidencias}\label{tab:casos-resumen}\\
\toprule
\textbf{ID} & \textbf{Nombre} & \textbf{Área} & \textbf{Tipo} & \textbf{Resultado} & \textbf{Observaciones} \\
\midrule
\endfirsthead
\toprule
\textbf{ID} & \textbf{Nombre} & \textbf{Área} & \textbf{Tipo} & \textbf{Resultado} & \textbf{Observaciones} \\
\midrule
\endhead
CP-001 & TODO & TODO & TODO & TODO & TODO \\
CP-002 & TODO & TODO & TODO & TODO & TODO \\
CP-003 & TODO & TODO & TODO & TODO & TODO \\
\bottomrule
\end{longtable}
}
\end{widepage}

\section{Hallazgos o defectos encontrados}

TODO: Presente los hallazgos, defectos, problemas, riesgos o incumplimientos detectados durante la evaluación.

\begin{widepage}
{\small
\setlength{\tabcolsep}{3pt}
\begin{longtable}{p{1.8cm} p{3.3cm} p{5.2cm} p{1.8cm} p{1.8cm} p{4.6cm} p{2cm} p{2cm}}
\caption{Registro formal de hallazgos o defectos}\label{tab:defectos}\\
\toprule
\textbf{ID} & \textbf{Área} & \textbf{Descripción} & \textbf{Severidad} & \textbf{Prioridad} & \textbf{Resultado esperado} & \textbf{Estado} & \textbf{Resp.} \\
\midrule
\endfirsthead
\toprule
\textbf{ID} & \textbf{Área} & \textbf{Descripción} & \textbf{Severidad} & \textbf{Prioridad} & \textbf{Resultado esperado} & \textbf{Estado} & \textbf{Resp.} \\
\midrule
\endhead
DEF-001 & TODO & TODO & TODO & TODO & TODO & TODO & TODO \\
DEF-002 & TODO & TODO & TODO & TODO & TODO & TODO & TODO \\
\bottomrule
\end{longtable}
}
\end{widepage}

\section{Análisis de resultados}

TODO: Interprete los resultados obtenidos. Relacione métricas, evidencias, requisitos, hallazgos y criterios de aceptación.

\section{Evaluación general}

\subsection{Cumplimiento funcional}

TODO: Analice el cumplimiento funcional observado.

\subsection{Calidad no funcional}

TODO: Analice rendimiento, accesibilidad, seguridad, usabilidad, mantenibilidad u otros atributos relevantes.

\subsection{Riesgos identificados}

\begin{table}[H]
\centering
\caption{Riesgos identificados}
\label{tab:riesgos}
\begin{tabularx}{\textwidth}{p{3cm} Y p{3cm} p{3cm}}
\toprule
\textbf{Riesgo} & \textbf{Descripción} & \textbf{Impacto} & \textbf{Mitigación} \\
\midrule
TODO & TODO & TODO & TODO \\
TODO & TODO & TODO & TODO \\
TODO & TODO & TODO & TODO \\
\bottomrule
\end{tabularx}
\end{table}

\section{Conclusiones}

\begin{enumerate}[label=\arabic*.]
    \item TODO: Conclusión 1.
    \item TODO: Conclusión 2.
    \item TODO: Conclusión 3.
\end{enumerate}

\section{Recomendaciones}

\begin{enumerate}[label=\arabic*.]
    \item TODO: Recomendación 1.
    \item TODO: Recomendación 2.
    \item TODO: Recomendación 3.
\end{enumerate}

\section{Anexos}

TODO: Incluya evidencias, reportes técnicos, criterios adicionales o material complementario.

\subsection{Anexo A: evidencias visuales}

\begin{longtable}{p{3cm} p{4.2cm} p{7cm}}
\caption{Evidencias visuales utilizadas}\label{tab:anexo-evidencias}\\
\toprule
\textbf{Área} & \textbf{Casos relacionados} & \textbf{Evidencia documentada} \\
\midrule
\endfirsthead
\toprule
\textbf{Área} & \textbf{Casos relacionados} & \textbf{Evidencia documentada} \\
\midrule
\endhead
TODO & TODO & TODO \\
TODO & TODO & TODO \\
TODO & TODO & TODO \\
\bottomrule
\end{longtable}

\subsection{Anexo B: reportes o respaldos técnicos}

\begin{table}[H]
\centering
\caption{Reportes técnicos considerados como respaldo}
\label{tab:anexo-reportes}
\begin{tabularx}{\textwidth}{p{3.8cm} p{4.2cm} Y}
\toprule
\textbf{Herramienta / Fuente} & \textbf{Reporte o resultado} & \textbf{Uso dentro del análisis} \\
\midrule
TODO & TODO & TODO \\
TODO & TODO & TODO \\
TODO & TODO & TODO \\
\bottomrule
\end{tabularx}
\end{table}

\subsection{Anexo C: criterios de clasificación}

\begin{table}[H]
\centering
\caption{Criterios de severidad aplicados a hallazgos}
\label{tab:anexo-severidad}
\begin{tabularx}{\textwidth}{p{3cm} Y}
\toprule
\textbf{Severidad} & \textbf{Criterio utilizado} \\
\midrule
Crítica & Bloquea un flujo principal completo, impide la operación esencial o expone información sensible. \\
Alta & Afecta una función importante y ofrece una alternativa limitada o poco clara. \\
Media & Afecta experiencia, consistencia o resultados parciales sin bloquear completamente el uso. \\
Baja & Corresponde a un problema visual, textual, de accesibilidad menor o de mejora de comunicación. \\
\bottomrule
\end{tabularx}
\end{table}

\begin{table}[H]
\centering
\caption{Criterios de prioridad aplicados a hallazgos}
\label{tab:anexo-prioridad}
\begin{tabularx}{\textwidth}{p{3cm} Y}
\toprule
\textbf{Prioridad} & \textbf{Criterio utilizado} \\
\midrule
Alta & Debe atenderse antes de una entrega porque afecta fuertemente el resultado o impide validar un requisito relevante. \\
Media & Debe documentarse y corregirse cuando se planifique mantenimiento del área afectada. \\
Baja & Puede gestionarse como mejora futura si no compromete el flujo principal ni la seguridad. \\
\bottomrule
\end{tabularx}
\end{table}

\end{document}
